\documentclass[11pt]{article}

% Packages
\usepackage[margin=1in]{geometry}
\usepackage{setspace}
\usepackage{titlesec}
\usepackage{graphicx}
\usepackage{booktabs}
\usepackage{amsmath}
\usepackage{hyperref}
\usepackage{float}
\usepackage{natbib}

\setstretch{1.15}

% Title formatting
\titleformat{\section}{\large\bfseries}{\thesection}{1em}{}
\titleformat{\subsection}{\normalsize\bfseries}{\thesubsection}{1em}{}

% Title info
\title{Airbnb NLP project}
\author{
Xianriu Cao \and Elvis Casco \and Román Feria
}


\date{\today}

\begin{document}

\maketitle

\begin{abstract}
This project uses natural language processing (NLP) to analyze Airbnb reviews in Barcelona in order to identify the factors that guests value most and to provide practical insights for real estate investors and Airbnb hosts. Drawing on review data collected over a specified period, the study examines how guests evaluate their experiences and which dimensions are most closely associated with satisfaction. The methodology combines TF--IDF, LDA topic modeling, and supervised machine learning. TF--IDF is used to represent the textual data, LDA to extract the main themes from the review corpus, and a classification model to predict which topics guests value the most. The findings suggest that guests focus primarily on host interaction, room quality, location and accessibility, and comfort and amenities, ranked in that order of importance. The project therefore offers a data-driven perspective on guest preferences in Barcelona's short-term rental market.
\end{abstract}

\section{Introduction}

\subsection{Background}
This project is motivated by the growing importance of the short-term rental market in contemporary urban economies. In Barcelona, tourism plays a central economic role, and platforms such as Airbnb have become a significant part of the accommodation sector and an increasingly relevant consideration for property owners and investors. At the same time, the expansion of real estate investment has intensified interest in understanding which property characteristics are most likely to generate positive guest experiences and sustain competitive performance in the market.

In this context, guest reviews represent a particularly valuable source of evidence. Unlike numerical ratings alone, textual reviews provide detailed and direct accounts of what visitors appreciated or criticized during their stay. They therefore make it possible to identify the factors that most strongly shape guest satisfaction, such as the quality of host interaction, room conditions, location, and overall comfort.

The analysis of reviews through natural language processing offers a systematic and data-driven way to uncover these patterns at scale. \citet{safari2025sentiment} classified the sentiment of reviews across six U.S.\ regions and found that more than 90\% of them are positive, and that what matters for host success is actually the polarity of reviews rather than how many reviews a listing has. Similarly, \citet{becerra2025reviews} showed that the sentiment expressed in reviews has a stronger effect on Airbnb prices than quantitative metrics like star ratings. \citet{perezsanchez2023sustainable} went further by building a price prediction model that combines machine learning with sentiment scores from review text.

This document seeks to contribute practical insights for Airbnb hosts and prospective investors by showing how guest preferences are expressed in reviews and which dimensions appear to matter most in Barcelona's short-term rental market.

\subsection{Main Idea}
The main aim of this project is to analyze Airbnb reviews from Barcelona in order to identify the principal dimensions of the guest experience reflected in customer feedback.

Barcelona is a particularly interesting case for studying Airbnb, and it has been used as a study site in a number of papers. With over 12 million international tourists visiting each year, the city faces intense pressure from short-term rentals, and this has made it one of the most regulated Airbnb markets in Europe. Recent text mining work on Airbnb data from Mediterranean cities, including Barcelona, has confirmed the richness and diversity of review content in these settings \citep{fernandezmorales2025mapping}.

On the regulatory side, Barcelona's city council launched PEUAT in 2017, a special urban plan that essentially froze the number of new short-term rental licenses in high-density zones \citep{ajuntament2017peuat}. However, enforcement has remained a challenge, and unlicensed listings continue to appear on the platform despite inspection efforts \citep{seguillinares2019monitoring}. This regulatory tension makes Barcelona an especially relevant setting for analyzing what guests and hosts discuss in reviews, as it may shape the kinds of experiences that are reported.

The data for this project comes from Inside Airbnb \citep{cox2016airbnb}, an open data initiative that scrapes and publishes listing and review data for cities around the world, including Barcelona.

Using natural language processing techniques, the study seeks to uncover the dominant themes that emerge from the review corpus and to assess their relative importance. We hypothesize that guest evaluations are structured around a limited number of recurring themes and that these themes do not contribute equally to the overall perception of a stay. More specifically, we expect interpersonal and experiential factors, such as host interaction and room quality, to be more prominent than supplementary aspects such as comfort and amenities. In this way, the project aims to provide a clearer understanding of what guests value most in Barcelona's short-term rental market.

\subsection{Motivation}
This topic is important from both a practical and a methodological perspective. From a practical point of view, understanding what guests value most in Airbnb accommodation can help hosts improve the quality of their service and enable prospective investors to make more informed decisions about which property characteristics are likely to matter in a competitive short-term rental market. In a city such as Barcelona, where tourism plays a major economic role, identifying the factors that shape guest evaluations is particularly relevant for understanding how short-term rentals create value.

From a methodological perspective, this project is motivated by the potential of natural language processing to extract meaningful patterns from large volumes of unstructured textual data. The analysis begins with TF--IDF, which is used to transform the review corpus into a document-term matrix and to represent the relative importance of words across the dataset. LDA is then applied to identify the main latent topics that structure guest reviews. Finally, these topics are analyzed in order to compare their relative prominence and to determine which dimensions of the guest experience appear most central in customer feedback. Rather than establishing causal relationships, this approach provides a systematic way to uncover and interpret the themes that guests emphasize most strongly in Barcelona's Airbnb market.

\section{Data}

\subsection{Data Sources}
The dataset used in this study was obtained from Inside Airbnb, an open-access platform that provides detailed information on Airbnb listings and reviews. This source is particularly well suited to our analysis because it offers large-scale, real-world data on guest experiences and on the characteristics of the short-term rental market in Barcelona.

\subsection{Exploratory Data Analysis (EDA)}
Before applying text mining techniques, we performed an exploratory analysis of the dataset in order to better understand its structure and general characteristics. Exploratory data analysis helps identify potential issues in the data, such as missing values, extremely short or long reviews, and irregularities in the textual content.

In this stage, we focused on examining the basic properties of the review corpus. This included inspecting the number of observations, exploring examples of the review text, and analyzing the distribution of review lengths. Studying the length of reviews can provide useful insights into how users typically describe their experiences on the platform and whether the dataset contains outliers that may affect the analysis.

Additionally, we analyzed the most frequent words appearing in the corpus. This provided an initial overview of the dominant vocabulary used by guests when describing their stays. Such insights can highlight commonly discussed aspects of Airbnb experiences, such as location, cleanliness, communication with the host, or the overall quality of the accommodation.

The results of this exploratory analysis helped guide the preprocessing and modeling steps that follow, ensuring that the text mining techniques were applied to a well-understood and properly prepared dataset.

The dataset consists of Airbnb reviews from Barcelona collected over the selected period of analysis. Because the reviews were originally written in multiple languages, the preprocessing stage began with translation into English in order to ensure linguistic consistency across the corpus. Variables not required for the study were then removed, and the final dataset was reduced to two essential fields: the review identifier and the review text.

One important finding is that the majority of reviews are relatively short. On the one hand, this suggests that guests often express their evaluations in a concise and selective manner, focusing mainly on the aspects of the stay they considered most memorable or most important. On the other hand, the brevity of many reviews also creates a methodological challenge for topic modeling, since LDA generally performs more effectively on longer documents that contain richer contextual information. To address this issue, the corpus was carefully preprocessed and the number of topics was restricted to a small and interpretable set, allowing the analysis to capture broad thematic patterns rather than highly specific distinctions.

In addition to review length, the EDA examined the frequency distribution of words and the most common terms in the corpus, thereby identifying the dominant lexical patterns that inform the later stages of the analysis.

% Example figure placeholder
\begin{figure}[H]
    \centering
    % \includegraphics[width=0.7\textwidth]{eda_plot.png}
    \caption{Example EDA figure.}
    \label{fig:eda}
\end{figure}

\section{Methodology}

\subsection{Sentiment Analysis}
Sentiment analysis has been used extensively in tourism research to make sense of large volumes of review text. Among the most popular tools for this is VADER, a lexicon-based model designed for social media text that assigns polarity scores to sentences \citep{hutto2014vader}. It is particularly attractive because it does not require labeled training data and works reasonably well out of the box.

One recurring issue in the Airbnb context is the so-called positivity bias: the overwhelming majority of reviews tend to be positive, which creates a skewed distribution. This is a problem for any classification model, since the negative class is severely underrepresented. Common strategies to deal with this include class weighting and oversampling techniques \citep{he2009learning}, both of which are used in practice across the literature.

\citet{santos2022expressing} looked at Airbnb reviews in Fortaleza, Brazil, and identified the main themes driving satisfaction and dissatisfaction, while \citet{agarwal2024investigating} took a more technical approach, combining NLP with Random Forest and GIS to segment listings by review-based sentiment, specifically trying to get around the positivity bias in numerical scores. Outside of Airbnb, it is also worth mentioning the work of \citet{kaya2024analysing}, who compared Booking.com and Expedia using sentiment analysis and topic modeling. They found, for instance, that Booking.com reviews focused more on ease of use while Expedia users complained more about cancellation and refund issues. This kind of cross-platform comparison is useful for understanding whether the patterns we observe in Airbnb reviews are platform-specific or more general.

\subsection{Overview of the Approach}
Our methodological framework combines TF--IDF, LDA, and supervised learning to analyze the main dimensions of guest evaluations in Airbnb reviews. TF--IDF is first used to construct a weighted document-term matrix, allowing the textual corpus to be represented numerically according to the relative importance of terms across reviews. LDA is then applied to identify the five most prominent latent topics in the corpus, thereby capturing the principal thematic dimensions of the guest experience. These outputs are subsequently incorporated into a supervised learning framework. Using Random Forest with cross-validation, we assess how the document-term features and topic structure contribute to the analysis, with the aim of determining which topics are most strongly associated with guest evaluations. Taken together, these methods provide a systematic NLP-based approach to identifying the aspects of the Airbnb experience that appear most important to guests.

\subsection{Preprocessing and Feature Extraction}
Before any analysis can be performed on review text, a series of preprocessing steps are needed to clean the data. These typically include lowercasing, removing punctuation and stopwords, and applying either stemming or lemmatization to reduce words to their base forms \citep{manning2008introduction}. In a multilingual setting like Barcelona, where reviews come in Spanish, Catalan, French, English, and several other languages, an additional step is needed: translating everything to a common language, usually English. While translation inevitably introduces some noise, neural machine translation tools like Google Translate have reached a quality level where this is generally acceptable for downstream NLP tasks.

\subsection{TF--IDF}
Once the text is clean, it needs to be converted into numerical features. TF--IDF (Term Frequency--Inverse Document Frequency) is still one of the most common representations for this purpose, despite being over fifty years old in concept \citep{sparckjones1972statistical, ramos2003using}. A useful extension is to include bigrams and trigrams in the vocabulary, which allows the model to capture multi-word expressions like ``check in,'' ``public transport,'' or ``well equipped'' --- phrases that carry important meaning in the Airbnb domain but would be lost if words were treated individually.

TF--IDF was used in this study for two main purposes. First, it transformed the review corpus into a weighted document-term matrix, providing the numerical representation required for subsequent modeling. Second, it helped identify the terms that were most distinctive within the corpus, thereby supporting interpretation of the lexical patterns in guest feedback.

TF--IDF assigns each term a weight based on its frequency within a document and its distribution across the corpus as a whole. Terms that appear frequently in a given review but relatively infrequently across other reviews receive higher weights, whereas terms that are common across the dataset receive lower weights. This makes TF--IDF useful for distinguishing general vocabulary from terms that are more informative or discriminative.

The choice of TF--IDF parameters was guided by exploratory analysis. A term-frequency distribution graph was first generated in order to examine how word frequencies were distributed across the dataset. Particular attention was paid to points at which the frequency bins displayed a marked jump, as these shifts suggested a meaningful threshold between overly common terms and more informative vocabulary. On this basis, we set \texttt{min\_df = 200}, retaining only terms that appeared in at least 200 documents, and \texttt{max\_df = 0.8}, excluding terms that appeared in more than 80\% of documents. This helped reduce noise from extremely rare words while preventing excessively frequent terms from dominating the representation.

An additional advantage of the TF--IDF vectorizer is that it allows us to examine the attribute \texttt{vectorizer.idf\_}, which returns the inverse document frequency (IDF) values for each term in the learned vocabulary. Formally, the IDF of a term \( t \) is defined as

\[
\mathrm{IDF}(t) = \log\left(\frac{N}{\mathrm{DF}(t)}\right) + 1
\]

where \( N \) is the total number of documents in the corpus, and \( \mathrm{DF}(t) \) is the number of documents in which term \( t \) appears.

The IDF measures how rare or informative a term is across the dataset. In our corpus, terms such as ``apartment'' and ``stay'' appear in a large share of reviews and therefore receive relatively low IDF values, indicating that they belong to the general vocabulary of Airbnb reviews rather than to any specific dimension of guest satisfaction. By contrast, less frequent terms tend to have higher IDF values and are therefore more useful for identifying distinctive aspects of the guest experience. Since \texttt{vectorizer.idf\_} captures only the global importance of a term across documents, it should be interpreted together with term frequency, as shown in the standard TF--IDF relationship:

\[
\mathrm{TF\text{-}IDF} = \mathrm{TF} \times \mathrm{IDF}.
\]

Overall, TF--IDF provided both the weighted textual representation required for later modeling and a useful basis for identifying the vocabulary that was most informative in the review corpus.

\subsection{LDA}
Latent Dirichlet Allocation, or LDA \citep{blei2003latent}, is probably the most widely used method for discovering topics in collections of documents in an unsupervised way. The basic idea is that each document is a mixture of topics, and each topic is characterized by a distribution over words. One of the main practical challenges with LDA is deciding how many topics to extract; there is no single correct answer: coherence scores \citep{roder2015exploring} provide a quantitative guide, but ultimately the topics also need to make sense to a human reader.

Several studies have applied LDA specifically to Airbnb reviews with interesting results. \citet{jang2020determinants} ran it on more than a million reviews from New York City and managed to extract 43 distinct topics, which they grouped into categories like location, unit characteristics, host quality, and general evaluations. Working with a smaller but still substantial dataset, \citet{chong2020comparative} compared reviews from Hong Kong (about 186,000) and Singapore (about 94,000), and found notable differences: 12 topics emerged for Hong Kong versus only 5 for Singapore, pointing to cross-cultural variation in what guests choose to write about. \citet{luo2022visual} proposed combining topic modeling with sentiment analysis to build a visual recommendation system for peer-to-peer accommodation.

What these studies have in common is that the topics tend to revolve around a few recurring dimensions: the physical property and its amenities, the host and communication, the location and its accessibility, and the overall value for money. This consistency across different cities and datasets suggests that the method is capturing something real about what guests care about.

We sought to determine which values for \texttt{min\_df} and \texttt{max\_df} would yield the highest coherence score. We therefore experimented with different values and different numbers of topics. The coherence score did not vary substantially across parameter combinations, so we selected a number of topics that produced the most interpretable groupings according to human judgment.

\subsection{Supervised Learning}
While lexicon-based tools like VADER are useful for quick sentiment scoring, supervised machine learning models trained on labeled data tend to be more accurate when enough training examples are available. The use of TF--IDF features with classifiers like Logistic Regression and Random Forest is a well-tested combination in NLP \citep{pang2002thumbs, liu2012sentiment}.

Logistic Regression is often the starting point because of its interpretability: the model coefficients tell you directly which words are most associated with positive or negative reviews, which can be very informative \citep{jurafsky2023speech}. Random Forest, on the other hand, can pick up on non-linear interactions between features and provides its own feature importance measure \citep{breiman2001random}. In practice, both models tend to perform reasonably well on review classification tasks, though Logistic Regression has the edge when it comes to understanding \textit{why} the model makes a particular prediction.

Given the class imbalance problem mentioned earlier, where positive reviews vastly outnumber negative ones, it is important to use techniques like class weighting (e.g., setting \texttt{class\_weight="balanced"}) so that the model does not simply learn to predict ``positive'' for every review \citep{he2009learning}.

\section{Results}
% TODO: Present the main findings of the analysis.

\section{Conclusion}
% TODO: Summarize the main findings, restate the contribution of the project, and provide a short takeaway.

\section{Limitations}
% TODO: Discuss the main limitations of the data, methods, and interpretation of the results. You may also mention possible future improvements.

\bibliographystyle{apalike}
\bibliography{references}

\end{document}
